(1) 角动量合成规则给 \(j=|l-s|,\ldots,l+s\)，因此 \(j=1/2\) 或 \(j=3/2\)。总磁量子数分别有 \(2j+1=2\) 与4个取值，恰与轨道三态乘自旋两态的六维空间相符。

(2) 利用
\[
2\boldsymbol L\cdot\boldsymbol S
=\boldsymbol J^2-\boldsymbol L^2-\boldsymbol S^2,
\]
对给定 j 的能移
\[
\Delta E_j=\frac\lambda2[j(j+1)-l(l+1)-s(s+1)].
\]
代入 l=1、s=1/2，得 \(\Delta E_{1/2}=-\lambda\)，\(\Delta E_{3/2}=\lambda/2\)。

(3) 总能级为 \(E_{1/2}=E_0-\lambda\)、\(E_{3/2}=E_0+\lambda/2\)，相邻间隔大小 \(3|\lambda|/2\)。若 \(\lambda >0\)，则 j=1/2 较低；若 \(\lambda <0\)，顺序反转。没有外场时每个 j 的 \(m_j\) 态仍简并。

(4) 对全部六个未偏置态按态数加权的平均能移
\[
\frac{2(-\lambda)+4(\lambda/2)}{6}=0.
\]
这也对应自旋轨道耦合在完整 l=1、s=1/2 张量积空间上的迹为零。
